% !PRINTED 223
\weberpage{223}
%<<CONT
deren Zahl gerade 60 betr\"agt. Dass sie alle von einander verschieden
 sind, sieht man, wenn man sie in der Form darstellt:

\weq{11}{\Theta^r = \begin{pmatrix} \varepsilon^r, & 0 \\ 0, & 1 \end{pmatrix},}

\weq{12}{\psi\Theta^r = \begin{pmatrix} 0, & 1 \\ -\varepsilon^r, & 0 \end{pmatrix},}

\weq{13}{\Theta^r\chi\Theta^s = \begin{pmatrix} \varepsilon^r\omega, & \varepsilon^{r-s}\beta \\ 1, & -\varepsilon^{-s}\omega \end{pmatrix},}

\weq{14}{\begin{aligned}
\Theta^r\chi\psi\Theta^s &= \begin{pmatrix} -\varepsilon^{r+s}\beta, & \varepsilon^r\omega \\ \varepsilon^s\omega, & 1 \end{pmatrix} \\
&= \begin{pmatrix} -\varepsilon^r\beta\omega^{-1}, & \varepsilon^{r-s} \\ 1, & \varepsilon^{-s}\omega^{-1} \end{pmatrix},
\end{aligned}}

von denen, da $-\omega^2$ keine Potenz von $\varepsilon$ ist, ersichtlich keine
zwei einander gleich sind.

Die Ikosa\"edergruppe hat, wie wir gesehen haben, 12 f\"unf-
z\"ahlige Pole. Je zwei dieser Pole sind die Pole von vier Substitutionen
 5\textsuperscript{ten} Grades, die mit der Identit\"at zusammen einen
f\"unfgliedrigen Cyklus bilden. Folglich giebt es in der Gruppe
24 Elemente 5\textsuperscript{ten} Grades. Ebenso giebt es 20 dreiz\"ahlige Pole,
die zu 20 Substitutionen dritter Ordnung f\"uhren, und 30 zwei-
z\"ahlige Pole, die zu zweien die Pole von je einer Substitution
zweiter Ordnung sind, so dass es 15 Substitutionen zweiter Ordnung
 giebt. Dies giebt mit der Identit\"at zusammen

\[
24 + 20 + 15 + 1 = 60.
\]

Die Bestimmung von $\beta$ in der Substitution $\chi$ ergiebt sich
nun durch Betrachtung der Grade von (13) und (14). Wir bilden
dazu zun\"achst f\"ur eine beliebige lineare Substitution $S = \begin{pmatrix} a, & b \\ c, & d \end{pmatrix}$
die zweite und dritte Wiederholung:

\[
\begin{aligned}
S^2 &= \begin{pmatrix} a^2 + bc, & b(a+d) \\ c(a+d), & bc + d^2 \end{pmatrix} \\
S^3 &= \begin{pmatrix} a^3 + 2abc + bcd, & b(a^2+bc+ad+d^2) \\ c(a^2+bc+ad+d^2), & d^3+abc+2bcd \end{pmatrix},
\end{aligned}
\]

und wenn wir von dem Falle absehen, dass $b$ oder $c$ verschwindet
, der hier nicht in Betracht kommt, so erhalten wir die
nothwendige und hinreichende Bedingung:
%CONT>>
